\section{Introduction}
\label{sec:introduction}
%======================================================================

The choice of square prefixes is not cosmetic.  It gives a cutoff with an exact square root, which turns a complicated smoothness condition into a clean one-large-prime decomposition.  The squared complex coordinate is then chosen because it records the product and factor imbalance at the same time.

Set
\[
 X_n=(n+1)^2-1,\qquad S_n=M(X_n).
\]
Why stop immediately before a square?  Because if $m<R^2$ is squarefree, it cannot contain two prime factors larger than $R$.  Every source is therefore either already $R$-smooth or has a unique representation
\[
 m=cq,\qquad c<R<q,
\]
with $q$ prime.  This gives the exact identity
\[
 S_n=A_n-T_n.
\]
The identity is elementary; the difficult question is why the two much larger terms $A_n$ and $T_n$ track each other so closely.

The geometry is designed to make that question visible.  For the factor pair $cq=m$, define
\[
 a=\frac{c+q}{2},\qquad b=\frac{q-c}{2}.
\]
Then
\[
 (a+bi)^2=m+i\frac{q^2-c^2}{2}.
\]
The real coordinate remembers the integer $m$, while the imaginary coordinate measures how uneven the factor pair is.  Fixed cofactors become parabolas.  The complete prefix is cut out by a vertical line.  Smoothness is cut out by a second moving parabola.  A source first enters the prefix, then spends time in the transport region, and finally becomes smooth when the moving boundary reaches its largest prime factor.

This viewpoint separates three logically different mechanisms:
\begin{enumerate}[leftmargin=1.7em,itemsep=0.35em]
\item \emph{exact bookkeeping}: the transport-to-smooth crossing cancels identically in $A-T$;
\item \emph{geometric dilution}: the quadratic map expands area by order $n^2$ at square scale $n$;
\item \emph{analytic cancellation}: the remaining high-height contribution must be controlled through one signed Gram form, including its off-diagonal terms.
\end{enumerate}
The first two are proved here.  The third is the central open estimate.  The coordinate change does not prove cancellation by itself; it identifies the correct state space, scale, and interaction that a proof must control.

The main exact results are summarized below.

\begin{theorem}[Exact square-root decomposition]
\label{thm:intro-decomposition}
For $R_n=n+1$ and $X_n=R_n^2-1$, define
\begin{align*}
 A_n&:=\sum_{\substack{m\le X_n\\\Pplus(m)\le R_n}}\mu(m),\\
 T_n&:=\sum_{c<R_n}\mu(c)
 \left[
 \pi\!\left(\left\lfloor\frac{X_n}{c}\right\rfloor\right)-\pi(R_n)
 \right].
\end{align*}
Then
\[
 \boxed{S_n=A_n-T_n.}
\]
\end{theorem}

\begin{theorem}[Squared-space recovery]
\label{thm:intro-recovery}
Let $w=x+iy$ be generated by a positive factor pair $cq=x$ through
\[
 y=\frac{q^2-c^2}{2}.
\]
With $\rho=|w|$,
\[
 \boxed{c^2=\rho-y,\qquad q^2=\rho+y.}
\]
Hence the squared point determines the ordered factor pair exactly.
\end{theorem}

\begin{theorem}[Cofactor-curve cutoff intersection]
\label{thm:intro-intersection}
For fixed $c>0$, put
\[
 \gamma_c(x)=\frac{x^2-c^4}{2c^2}.
\]
At cutoff $R>0$, put
\[
 \Gamma_R(x)=\frac{R^4-x^2}{2R^2}.
\]
Then
\[
 \boxed{\gamma_c(x)=\Gamma_R(x)\iff x=cR.}
\]
The curves meet orthogonally at their intersection.  At an arithmetic point $x=cq$, the sign of $\gamma_c(cq)-\Gamma_R(cq)$ is the sign of $q-R$.
\end{theorem}

\begin{theorem}[Exact channel telescoping]
\label{thm:intro-telescoping}
As $R$ increases, a canonical point enters the prefix at $R=\sqrt{cq}$ and crosses from transport to smooth at $R=q$.  The signed mass crossing the moving parabola appears with opposite signs in the increments of $A$ and $-T$.  Therefore all internal channel transfers cancel from the increment of $A-T$; only points entering the new square strip remain.
\end{theorem}

\begin{theorem}[Low-height clustering and lifetime]
\label{thm:intro-low}
Let $m=cq\in B_n$ be canonical and let
\[
 Y_m=\frac{q^2-c^2}{2}.
\]
Then, for $H\ge0$,
\[
 \#\{m\in B_n:|Y_m|\le H\}
 \le\left\lfloor\frac{H}{n}\right\rfloor,
\]
and the count is zero when $H=n$.  For odd sources the bound improves to
$\lfloor H/(2n)\rfloor$.
If $Y_m>0$, the transport lifetime satisfies
\[
 L(m)\le\frac{2Y_m}{q}.
\]
Consequently, for fixed $\Lambda$, the canonical points with $|Y_m|\le\Lambda n$ have uniformly bounded block occupancy, and the transport members have uniformly bounded lifetime.
\end{theorem}

\begin{theorem}[Elementary canonical pointwise criterion]
\label{thm:intro-canonical-HS}
Fix $\Lambda>0$.  Extend the harmless conventions $\Pplus(1)=1$, $c_1=1$, and $Y_1=0$.  For the square block $B_j=[j^2,(j+1)^2)$ define
\[
 d_j^{\mathrm{high}}(\Lambda)
 :=
 \sum_{\substack{m\in B_j\\|Y_m|>\Lambda j}}\mu(m),
 \qquad
 S_N^{\mathrm{high}}(\Lambda)
 :=\sum_{0\le j\le N}d_j^{\mathrm{high}}(\Lambda).
\]
Then the following are equivalent:
\begin{enumerate}[label=\textup{(\roman*)},leftmargin=2.2em,itemsep=0.25em]
\item the Riemann Hypothesis;
\item for every $\eps>0$,
\begin{equation}
 \boxed{
 |S_N^{\mathrm{high}}(\Lambda)|^2
 \ll_{\eps,\Lambda}N^{2+\eps};
 }
 \tag{PHS}
 \label{eq:intro-canonical-PHS}
\end{equation}
\item for every $\eps>0$ and $1\le H\le N$,
\begin{equation}
 \sum_{n=N}^{N+H-1}|S_n^{\mathrm{high}}(\Lambda)|^2
 \ll_{\eps,\Lambda}HN^{2+\eps}.
 \tag{HS}
 \label{eq:intro-canonical-HS}
\end{equation}
\end{enumerate}
The equivalence uses only the exact low/high recombination, the elementary low-sector bound, square-endpoint interpolation, and the classical Mertens criterion.  The implication \textup{(ii)}$\Rightarrow$\textup{(iii)} follows from $N+h<2N$; the reverse implication is the single choice $H=1$.
\end{theorem}

\begin{theorem}[Quadratic conformal expansion]
\label{thm:intro-jacobian}
The map
\[
 \Phi(c,q)=\left(cq,\frac{q^2-c^2}{2}\right)
\]
has
\[
 D\Phi(c,q)^{\mathsf T}D\Phi(c,q)
 =(c^2+q^2)I,
\]
and
\[
 \det D\Phi(c,q)=c^2+q^2\ge2cq.
\]
At square scale $cq\asymp n^2$, lengths expand by at least order $n$ and areas by at least order $n^2$.
\end{theorem}

\begin{theorem}[Exact dyadic cell compression]
\label{thm:intro-dyadic-cell}
For every $k\ge0$,
\[
 \boxed{
 \bigl(\mu(4k+1),\mu(4k+2),\mu(4k+3),\mu(4k+4)\bigr)
 =\bigl(\mu(4k+1),-\mu(2k+1),\mu(4k+3),0\bigr).
 }
\]
Consequently the complete-cell scalar is
\[
 \boxed{C_k=\mu(4k+1)-\mu(2k+1)+\mu(4k+3).}
\]
Among the three active slots $4k+1,4k+2,4k+3$, exactly one is divisible by $3$; the affected slot is $4k+3,4k+2,4k+1$ according as $k\equiv0,1,2\pmod3$.
\end{theorem}

\begin{theorem}[Exact small-modulus quadratic resonance]
\label{thm:intro-prime-square-24}
For every prime $q>3$,
\[
 \boxed{q^2\equiv1\pmod{24}.}
\]
Hence, for every integer $m$ and every prime interval containing no $2$ or $3$,
\[
 \boxed{
 \sum_{Q<q\le2Q\atop q\ {\rm prime}}
 e\!\left(\frac{m q^2}{24}\right)
 =e\!\left(\frac m{24}\right)
 \bigl(\pi(2Q)-\pi(Q)\bigr).
 }
\]
Equivalently, the quadratic prime phase is perfectly coherent at every
$\alpha=m/12$.  More generally, for $\alpha=a/r$, the rational phase
$e(aq^2/(2r))$ is periodic modulo $2r$; the correct reduced-residue model is therefore formed over $(\mathbb Z/2r\mathbb Z)^\times$.
\end{theorem}

The remaining theorem is the pointwise inequality \eqref{eq:intro-canonical-PHS}.  The local estimate \eqref{eq:intro-canonical-HS} is an exact equivalent.  No claim is made that the coordinate change alone supplies cancellation; it identifies the native arithmetic quantity, the correct scale, and the signed interaction that must be controlled.

%======================================================================
